Serial pipelines that reduce a per-element signal with a small encoder and
then apply a wide downstream model are the default in spectroscopic imaging,
and the same structure appears wherever a compact encoder feeds a large head.
We separate three ingredients of selective encoder adaptation in such models:
block curvature, residual excitation, and optimizer-dependent displacement.
Conditional bounds identify when a fast downstream module limits encoder
movement: a width-linear curvature disparity under an explicit head
hypothesis, a finite-time displacement bound under a residual envelope, a
sharp finite-horizon attribution bound, a transparent affine comparison of
joint and frozen training, and a scoped nonlinear existence result. An
exactly solvable serial cross-entropy instance exposes the role of the induced
training metric at a matched fitting threshold; there the displacement bound
is within a factor of six of the truth. Measurements on a production
spectral--spatial segmentation model show why scalar curvature can mislead:
rank-one inputs invert it, its width trend is carried by batch-normalization
gains, and a shallow-model trajectory experiment shows the residual leaving
the fast directions. Matched retraining comparisons show that a previously
reported freezing advantage cannot be attributed to architecture from that
protocol.
